% !TeX root = ../动手学强化学习笔记.tex
% 由主文档通过 \input 载入；请编译主文档。
\setcounter{chapter}{21} % 保留原书章号。
\chapter{模仿学习与离线强化学习}
\label{ch:rl-offline}

在线算法可以尝试新动作以纠正错误，固定数据下的学习则受到数据覆盖的根本限制。本章按可用信息与学习对象区分三个专题：模仿学习利用示范学习策略，其中行为克隆使用给定示范，交互式数据聚合还允许执行当前策略并继续查询专家；离线强化学习（offline reinforcement learning）使用固定的转移与奖励数据优化策略，训练期间不追加环境交互\cite{levine-offline}；逆强化学习与偏好反馈则从示范或比较信息学习奖励，所获奖励如何用于后续策略优化是另一层问题。因而，固定数据设定只约束其中相应的方法，不能概括全部模仿学习与奖励学习。

本章以分布偏移、数据覆盖和奖励可辨认性为主线，介绍行为克隆、保守价值估计与策略约束的代表机制，并概述离线评价的条件。逆强化学习与偏好反馈作为奖励学习的延伸，仅讨论问题设定、可辨认性及评价边界；本章不展开完整的奖励学习与策略优化训练流程。学习时应明确数据能否追加、监督信号是什么，以及当前估计的是策略、价值还是奖励。

\section{行为克隆与分布偏移}

模仿学习（imitation learning）从专家行为学习策略。行为克隆（behavior cloning）将示范集 $\mathcal D_E=\{(s_i,a_i)\}$ 看作监督数据，随机策略可最小化
\[
L_{\rm BC}(\theta)=-\frac1{|\mathcal D_E|}\sum_{(s,a)\in\mathcal D_E}\log\pi_\theta(a\mid s).
\]
它不需要奖励标签，却只在专家访问状态上直接约束行为。策略执行时一次错误可能把系统带到示范中没有出现的区域，后续错误因此增多；较低的示范分类误差不自动意味着较高的闭环任务收益。

若训练期间可以继续查询专家，可让当前策略执行并收集其实际访问状态，再请专家给动作标签，将新数据加入训练集。这种数据聚合针对分布偏移，但使用了额外专家查询，不能与只允许固定示范的方法按相同数据条件比较。

\section{固定数据下的价值学习与不可辨认性}

离线数据通常包含 $(s,a,r,s',d)$。在数据上做 Bellman 回归时，目标可能出现 $\max_{a'}Q(s',a')$，而某些后继动作从未被数据采集策略尝试。它们的高估既缺少直接监督，又会通过自举进入其他表项或网络参数。

\begin{example}[无限数据也未必足够]
一步任务有动作 $a,b$。数据策略始终选择 $a$，其奖励恒为 $1$。两个环境在 $a$ 上完全相同，但环境一的 $b$ 奖励为 $0$，环境二的 $b$ 奖励为 $2$。无论收集多少条这种数据，其观测分布都相同，而两个环境的最优动作不同。因此任何只读取这些数据的方法，都无法保证同时识别两者的最优策略。
\end{example}

这个例子是信息限制，不是模型容量不足。若额外假定动作间有已知结构，可以借结构推断未见动作，但结论便依赖该假设。离策略学习允许目标与行为不同，离线学习则禁止追加交互，两者应分别陈述。

\section{保守估计、策略约束与离线评价}

一种应对方式限制新策略偏离估计行为策略 $\widehat b$，例如在最大化估计回报时加入 $\mathbb E_sD_{\rm KL}(\pi(\cdot\mid s)\|\widehat b(\cdot\mid s))$ 惩罚；另一种方式降低缺少支持动作的价值。离散动作下，可考虑正则项
\[
\mathbb E_{s\sim\mathcal D}\left[\log\sum_a e^{Q(s,a)}-\mathbb E_{a\sim\mathcal D(\cdot\mid s)}Q(s,a)\right].
\]
这类正则可用于保守 Q 学习（conservative Q-learning，CQL）\cite{kumar-cql}。第一项倾向压低整体较高预测，第二项相对支持数据动作。它说明保守价值学习的一种机制，但权重、尺度及适用算法仍需具体定义；加入正则不能从没有信息的数据中恢复任意环境的真实最优价值。

离线策略评价（off-policy evaluation，OPE）可以使用重要性采样、模型估计或拟合价值等方法。重要性采样需要行为概率与覆盖，模型方法依赖预测准确性，价值方法依赖函数近似和数据支持。用同一批评价数据在许多候选策略中选最大估计者，会偏向正噪声，应保留独立评价数据或显式控制选择偏差。

\section{逆强化学习与偏好反馈}

逆强化学习（inverse reinforcement learning）从示范推断解释行为的奖励。奖励通常不可唯一辨认：第\ref{ch:rl-exploration}章的势函数塑形就给出不同奖励但相同最优策略的情况。恢复某一种奖励需要结构、先验或其他观测条件；能拟合示范不等于恢复了专家唯一的真实目标。

偏好学习改为观察两条轨迹的比较。若奖励模型给轨迹评分 $U_\psi(\tau)$，一种概率模型为 $\Pr(\tau_1\succ\tau_2)=\sigma(U_\psi(\tau_1)-U_\psi(\tau_2))$，可用比较标签拟合参数。该模型是一项统计假设；之后优化学得评分会改变策略访问的轨迹分布，可能进入比较数据稀少的区域。因此偏好预测、奖励模型和最终策略应分别评价。

\section{习题与实验}
\begin{enumerate}
\item 对不可辨认性例子说明：如果只能输出随机策略，能否在两个环境中同时保证最优？给出理由。
\item 构造示范状态上动作准确率很高、但一次偏离后回报很差的两阶段任务，比较是否允许专家再标注的两种数据条件。
\item 固定训练数据，比较行为克隆、直接价值学习和一种支持约束方案。按数据质量与覆盖分组报告结果，训练阶段不得额外访问评价环境。
\end{enumerate}





